\documentclass[11pt,a4paper]{article}

\usepackage[utf8]{inputenc}
\usepackage[T1]{fontenc}
\usepackage{amsmath,amssymb,amsthm}
\usepackage{hyperref}
\usepackage{geometry}
\usepackage{xcolor}
\usepackage{enumitem}
\usepackage{booktabs}
\usepackage{cleveref}

\geometry{margin=1in}

\newtheorem{theorem}{Theorem}[section]
\newtheorem{lemma}[theorem]{Lemma}
\newtheorem{definition}[theorem]{Definition}
\theoremstyle{remark}
\newtheorem{remark}[theorem]{Remark}

\title{\bfseries Lean 4 Infrastructure Toward the Analytic Prime Number Theorem\\
{\large Zeta-Function Formalization, Zero-Free-Region Targets, and Gap Analysis}}

\author{Edward Deepseek\\
  \texttt{Maxminicherry@gmail.com}
}

\date{\today}

\begin{document}

\maketitle

\begin{abstract}
We report on a buildable Lean~4 project developing infrastructure toward
formalizing de~la~Vall\'ee~Poussin's analytic proof of the Prime Number
Theorem (PNT).  The repository works directly with Mathlib's
\texttt{riemannZeta} API and contains formal statements and supporting lemmas
for the 3-4-1 trigonometric method, a compact zero-free region, quantitative
zero-free-region targets, pole-behavior facts, Hardy's theorem targets, and
PNT-equivalence statements.

The present development should be read as a verified framework with explicit
unproved target statements, not as a completed PNT formalization.  With
Lean~4.29.1 and Mathlib~4.29.1 the project builds successfully and contains no
syntactic \texttt{sorry} occurrences.  The remaining targets include Hardy's
integral asymptotics, the quantitative zero-free regions, the von Mangoldt
explicit formula, and standard RH/PNT error-term equivalences.
We give a precise
inventory of what is verified, what remains conjectural within the code, and
which Mathlib components would be needed to close the gaps.
\end{abstract}

%=====================================================================
\section{Introduction}
%=====================================================================

The Prime Number Theorem, stating that the number of primes up to $x$
satisfies $\pi(x) \sim x/\log x$, was proved independently by Hadamard and
de~la~Vall\'ee~Poussin in 1896. Their proofs established that the Riemann
zeta function $\zeta(s)$ has no zeros on the line $\Re(s)=1$, from which
the PNT follows via contour integration (Perron's formula).

Formalizing these proofs in an interactive theorem prover has been a
long-standing challenge. Previous formalizations have taken different routes:

\begin{itemize}[noitemsep]
  \item Avigad et al.\ (2007) formalized the elementary Selberg--Erd\H{o}s
    proof in Isabelle.
  \item Harrison (2009) formalized Newman's short analytic proof in HOL Light.
  \item The \texttt{PrimeNumberTheoremAnd} project (Kontorovich \& Tao) is
    formalizing the Wiener--Ikehara proof in Lean~4.
\end{itemize}

Notably absent, however, has been a complete formalization of the
\emph{original} analytic proof via the 3-4-1 inequality.  This project is an
intermediate step toward such a formalization in Lean~4, built on Mathlib.

Our contributions are:

\begin{enumerate}[noitemsep]
  \item A Lean~4 development of the trigonometric identity underlying
    de~la~Vall\'ee~Poussin's 3-4-1 method, together with a verified
    logarithmic-derivative series combination.
  \item A verified compact zero-free region for bounded height, plus target
    statements for quantitative zero-free regions of $\zeta(s)$ near
    $\Re(s)=1$.
  \item Supporting facts around zeta non-vanishing, pole-behavior wrappers,
    Euler products, zeta values, Hardy's theorem, and PNT-equivalence
    statements.
  \item A reproducible Lean~4.29.1 build and an explicit gap table listing the
    remaining unproved target statements.
\end{enumerate}

%=====================================================================
\section{Background}
%=====================================================================

\subsection{The Riemann Zeta Function}

For $\Re(s)>1$, the Riemann zeta function is defined by the Dirichlet series
$\zeta(s)=\sum_{n=1}^{\infty} n^{-s}$. It extends meromorphically to the
whole complex plane with a single simple pole at $s=1$ of residue~1. The
function satisfies the functional equation
\[
\zeta(1-s) = 2(2\pi)^{-s}\,\Gamma(s)\cos(\pi s/2)\,\zeta(s),
\]
which implies that $\zeta(-2n)=0$ for $n\in\mathbb{N}$ (the trivial zeros)
and that all other zeros lie in the critical strip $0<\Re(s)<1$.

\subsection{The von Mangoldt Function and Chebyshev Functions}

The von Mangoldt function $\Lambda(n)$ equals $\log p$ if $n=p^k$ is a prime
power, and $0$ otherwise. The Chebyshev $\psi$-function is
$\psi(x)=\sum_{n\leq x}\Lambda(n)$, and the PNT is equivalent to
$\psi(x)\sim x$.

A fundamental identity links $\Lambda$ with the logarithmic derivative of
$\zeta$:
\[
-\frac{\zeta'(s)}{\zeta(s)} = \sum_{n=1}^{\infty} \frac{\Lambda(n)}{n^{s}},
\qquad \Re(s)>1.
\]

\subsection{The 3-4-1 Inequality}

De~la~Vall\'ee~Poussin's key insight was the trigonometric inequality
\[
3 + 4\cos\theta + \cos 2\theta = 2(1+\cos\theta)^{2} \geq 0.
\]
For $\sigma>1$ and real $t$, using $\Re(n^{-s}) = n^{-\sigma}\cos(t\log n)$,
this yields
\[
3\,\Re\!\left(-\frac{\zeta'}{\zeta}(\sigma)\right)
+ 4\,\Re\!\left(-\frac{\zeta'}{\zeta}(\sigma+it)\right)
+ \Re\!\left(-\frac{\zeta'}{\zeta}(\sigma+2it)\right) \geq 0.
\]
From this, one deduces that $\zeta(1+it)\neq 0$ for all real $t$, and more
generally obtains zero-free regions.

%=====================================================================
\section{Formalization Architecture}
%=====================================================================

The formalization consists of seven Lean source files, structured to separate
concerns and minimize dependencies:

\begin{center}
\begin{tabular}{lrl}
\toprule
\textbf{File} & \textbf{Key Contents} & \textbf{Status} \\
\midrule
\texttt{RiemannPNT.lean} &
  Build entry point importing all top-level modules &
  \texttt{sorry}-free \\
\texttt{RiemannExplorer.lean} &
  RH statement, zeta definitions, functional equation, trivial zeros &
  \texttt{sorry}-free \\
\texttt{EulerAndLfunctions.lean} &
  Thin wrappers around Mathlib zeta, Euler-product, and L-function facts &
  \texttt{sorry}-free \\
\texttt{GammaResidue.lean} &
  $\operatorname{Res}(\Gamma,-n)=(-1)^{n}/n!$ &
  \texttt{sorry}-free \\
\texttt{HardyTheorem.lean} &
  Hardy Z-function framework; corrected targets for the missing moment estimates &
  \texttt{sorry}-free; targets unproved \\
\texttt{PrimeNumberTheorem.lean} &
  PNT forms, equivalences, $\operatorname{Li}(x)$ asymptotics, zero symmetry,
  bounded-height zero finiteness, explicit formula target &
  \texttt{sorry}-free; targets unproved \\
\texttt{ZeroFreeRegion.lean} &
  3-4-1 inequality, compact zero-free region, log derivative series, residue
  bounds &
  \texttt{sorry}-free; quantitative targets unproved \\
\bottomrule
\end{tabular}
\end{center}

\subsection{Verified Local Contributions}

The main project-local verified contributions are concentrated in
\texttt{ZeroFreeRegion.lean}.  They are the pieces of classical zeta-function
infrastructure needed before one can attempt the quantitative
de~la~Vall\'ee~Poussin zero-free region.

\paragraph{Real part of the logarithmic derivative.}
For $s=\sigma+it$ with $\sigma>1$, we prove
\[
\Re\!\left(-\frac{\zeta'}{\zeta}(s)\right)
  = \sum_{n=1}^{\infty}
      \frac{\Lambda(n)\cos(t\log n)}{n^\sigma}.
\]
The proof starts from Mathlib's complex L-series identity for the von
Mangoldt function and then extracts real parts term by term.  The key local
calculation is
\[
\Re(n^{-s}) = n^{-\sigma}\cos(t\log n),
\]
formalized by expanding $n^{-s}$ as $\exp(-s\log n)$ and using the fact that
the complex logarithm of a positive real natural number agrees with its real
logarithm.

\paragraph{The 3-4-1 combination.}
Using the preceding series representation, we prove
\[
3\Re\!\left(-\frac{\zeta'}{\zeta}(\sigma)\right)
+4\Re\!\left(-\frac{\zeta'}{\zeta}(\sigma+it)\right)
+\Re\!\left(-\frac{\zeta'}{\zeta}(\sigma+2it)\right)\ge 0.
\]
After expanding the three terms, the coefficient of each summand is
\[
3\cos(0\cdot\log n)+4\cos(t\log n)+\cos(2t\log n),
\]
which reduces to $3+4\cos(t\log n)+\cos(2t\log n)$.  The latter is
$2(1+\cos(t\log n))^2$, hence nonnegative.  The Lean proof also supplies the
summability hypotheses needed to combine the three infinite sums into one
\texttt{tsum}.

\paragraph{Compact zero-free strip.}
For every $T\ge2$, we prove that there is $d>0$ such that
\[
|\Im(s)|\le T,\quad \Re(s)\ge 1-d
  \quad\Longrightarrow\quad \zeta(s)\ne0.
\]
The checked proof uses an open-thickening compactness argument.  First,
\(\zeta\) is shown to be nonzero in a full neighborhood of \(1\), using the
residue statement on the punctured neighborhood and Mathlib's value at \(1\).
Next, the nonvanishing locus is proved open.  The compact vertical segment
\(\{s:\Re(s)=1,\ |\Im(s)|\le T\}\) is contained in this open set, so
\texttt{IsCompact.exists\_cthickening\_subset\_open} gives a uniform positive
thickness.  Points with $\Re(s)\ge1$ are handled directly by Mathlib's
nonvanishing theorem; points with $1-d\le\Re(s)<1$ are projected vertically to
the compact segment.

\paragraph{Residue-scale bounds on the real axis.}
For real $\sigma>1$, we prove
\[
1 < (\sigma-1)\Re(\zeta(\sigma)) \le \sigma.
\]
Since $\zeta(\sigma)$ is real on this interval, this is the usual
residue-scale estimate \(1<(\sigma-1)\zeta(\sigma)\le\sigma\).  The proof
uses Mathlib's \texttt{ZetaAsymptotics.zeta\_limit\_aux1}, comparing the
auxiliary nonnegative integral sum at $\sigma$ with its value at $1$ and using
the lower bound $\gamma>1/2$ for the Euler--Mascheroni constant.

\begin{center}
\begin{tabular}{lll}
\toprule
\textbf{Input already in Mathlib} & \textbf{Project-local result} & \textbf{Next missing step} \\
\midrule
Von Mangoldt L-series identity &
Real-part series for $-\zeta'/\zeta$ &
Quantitative log-derivative estimates \\
$\zeta(s)\ne0$ for $\Re(s)\ge1$ &
Compact zero-free strip at bounded height &
Uniform width $c/\log|t|$ \\
Residue of $\zeta$ at $1$ &
Bounds on $(\sigma-1)\zeta(\sigma)$ &
Growth estimates near the pole \\
\bottomrule
\end{tabular}
\end{center}

The dependency pattern toward the classical zero-free region is:
\[
\text{log-derivative series}
\longrightarrow
\text{3-4-1 inequality}
\longrightarrow
\text{lower bound for }-\zeta'/\zeta
\quad\text{combined with}\quad
\text{residue-scale bounds}.
\]
The remaining analytic input is a global growth/logarithmic-derivative bound,
usually obtained through Hadamard factorization or a Borel--Carath\'eodory
route.

\subsection{Prime Number Theorem Equivalences}

In \texttt{PrimeNumberTheorem.lean}, we define the three classical
formulations:

\begin{itemize}[noitemsep]
  \item \texttt{PNTForm1}: $\displaystyle\lim_{x\to\infty}
    \frac{\pi(x)\log x}{x}=1$
  \item \texttt{PNTForm2}: $\displaystyle\lim_{x\to\infty}
    \frac{\pi(x)}{\operatorname{Li}(x)}=1$
  \item \texttt{PNTForm3}: $\displaystyle\lim_{x\to\infty}
    \frac{\psi(x)}{x}=1$
\end{itemize}

We prove that all three are equivalent (\texttt{pnt\_forms\_equivalent}).
The proof of \texttt{PNTForm1~$\leftrightarrow$~PNTForm2} uses the
asymptotic $\operatorname{Li}(x)\sim x/\log x$, which we prove in full
detail by splitting the integral at $\sqrt{x}$ and bounding each part.
The equivalence \texttt{PNTForm2~$\leftrightarrow$~PNTForm3} uses the
Chebyshev bounds $|\psi(x)-\theta(x)|\leq 2\sqrt{x}\log x$ and the
asymptotic relationship between $\theta(x)$ and $\pi(x)$.

\subsection{The 3-4-1 Inequality}

The core lemma \texttt{trig\_identity\_nonneg} proves
$3+4\cos\theta+\cos 2\theta = 2(1+\cos\theta)^{2} \geq 0$ in Lean.

The main theorem \texttt{log\_deriv\_zeta\_nonneg\_combination} establishes
\[
3\cdot\Re\!\left(-\frac{\zeta'}{\zeta}(\sigma)\right)
+ 4\cdot\Re\!\left(-\frac{\zeta'}{\zeta}(\sigma+it)\right)
+ \Re\!\left(-\frac{\zeta'}{\zeta}(\sigma+2it)\right) \geq 0.
\]

The proof proceeds by:
\begin{enumerate}[noitemsep]
  \item Expressing each term as a Dirichlet series in $\Lambda(n)$ via
    Mathlib's \texttt{LSeries\_vonMangoldt\_eq\_deriv\_riemannZeta\_div}.
  \item Extracting the real part using
    $\Re(n^{-s}) = n^{-\sigma}\cos(t\log n)$.
  \item Combining the three series into
    $\sum \frac{\Lambda(n)}{n^{\sigma}}(3+4\cos(t\log n)+\cos(2t\log n))$.
  \item Applying \texttt{trig\_identity\_nonneg} termwise.
\end{enumerate}

\subsection{Compact Zero-Free Region}

The theorem \texttt{classical\_zero\_free\_region\_compact} states: for any
$T\geq 2$, there exists $d>0$ such that $\zeta(s)\neq 0$ whenever
$|\Im(s)|\leq T$ and $\Re(s)\geq 1-d$.

The proof uses only point-set topology: since $\{s:|\Im(s)|\leq T,\;
\Re(s)=1\}$ is compact and is contained in the open nonvanishing locus of
$\zeta$, there exists a uniform thickening on which $\zeta$ remains nonzero.
This compact result is independent of the 3-4-1 argument; the latter is needed
for the later quantitative width $c/\log|t|$.

\subsection{Pole Structure at \texorpdfstring{$s=1$}{s=1}}

We formalize the following bounds for real $\sigma>1$:
\[
1 < (\sigma-1)\,\zeta(\sigma) \leq \sigma,
\]
proved via the integral representation and the Euler--Mascheroni constant
bounds. This confirms the residue of $\zeta$ at $s=1$ is~1.

\subsection{Zero Symmetry and Trivial Zeros}

Using the functional equation, we prove:
\begin{itemize}[noitemsep]
  \item \texttt{riemannZeta\_ne\_zero\_of\_re\_le\_zero}: For $\Re(s)\leq 0$,
    the only zeros are at $s=-2,-4,-6,\dots$
  \item \texttt{nontrivial\_zero\_symmetric}: If $\rho$ is a non-trivial zero,
    then $1-\rho$ is also a non-trivial zero.
  \item \texttt{nontrivial\_zero\_in\_critical\_strip}: All non-trivial zeros
    satisfy $0<\Re(s)<1$.
\end{itemize}

%=====================================================================
\section{Infrastructure Gaps}
%=====================================================================

The current checkout contains no syntactic \texttt{sorry} occurrences.  The
unproved deep results are recorded as \texttt{Prop} target statements rather
than exported theorems.  Some are blocked by missing Mathlib infrastructure;
others are project-local proof obligations that still need to be completed.
Table~\ref{tab:gaps} summarizes the main blockers.

\begin{table}[ht]
\centering
\caption{Main unproved target statements and blockers.}
\label{tab:gaps}
\begin{tabular}{lll}
\toprule
\textbf{Theorem} & \textbf{Missing Component} & \textbf{Difficulty} \\
\midrule
\texttt{explicit\_formula\_von\_mangoldt}
  & Contour integration \& residue theorem & High \\
\texttt{rh\_iff\_optimal\_error}
  & Standard RH/PNT error-term theory & High \\
\texttt{classical\_zero\_free\_region}
  & Hadamard factorization \emph{or} Borel--Carath\'eodory & Medium \\
\texttt{vinogradov\_korobov\_zero\_free\_region}
  & Exponential sum estimates & Very High \\
Hardy's integral asymptotics
  & Corrected moment and special-function estimates & Medium--High \\
\bottomrule
\end{tabular}
\end{table}

\subsection{The Borel--Carath\'eodory Shortcut}

The quantitative zero-free region $\sigma \geq 1-c/\log|t|$ typically
requires Hadamard factorization theory. However, there is a lighter route
via the Borel--Carath\'eodory theorem (which bounds an analytic function by
its real part). Adding Borel--Carath\'eodory to Mathlib would remove one
major infrastructure obstacle, but the present repository would still need
additional zeta growth and logarithmic-derivative estimates before the full
classical PNT proof could be closed.

\subsection{The Explicit Formula}

The von Mangoldt explicit formula
\[
\psi(x) = x - \sum_{\rho} \frac{x^{\rho}}{\rho}
  - \frac{\zeta'(0)}{\zeta(0)} - \frac{1}{2}\log(1-x^{-2})
\]
requires contour integration and the residue theorem---both absent from
Mathlib as of writing. This is a larger undertaking but would unlock many
results beyond the PNT.

%=====================================================================
\section{Related Work}
%=====================================================================

\paragraph{Elementary proof in Isabelle.}
Avigad, Donnelly, Gray, and Raff (2007)~\cite{avigad2007pnt} formalized the
elementary Selberg--Erd\H{o}s proof in Isabelle/HOL. This remains the only
complete formalization of the PNT, but the elementary route avoids all
complex analysis and does not give error terms.

\paragraph{Newman's proof in HOL Light.}
Harrison (2009)~\cite{harrison2009pnt} formalized Newman's short analytic
proof, which uses contour integration over a rectangle. This work was
tighter in scope and did not develop the full zeta function infrastructure.

\paragraph{PrimeNumberTheoremAnd in Lean~4.}
Kontorovich and Tao's ongoing project~\cite{pntand} formalizes the PNT
via the Wiener--Ikehara theorem, a Fourier-analytic approach that avoids
much of the classical zeta theory.

\paragraph{Zeta function in Mathlib.}
Loeffler and Stoll have built the foundational zeta function theory in
Mathlib, including the functional equation, Euler product, and Dirichlet
$L$-functions~\cite{mathlib}. Our work directly extends this foundation.

\paragraph{Our contribution.}
Our formalization is the first to tackle the \emph{classical} analytic proof
directly, developing the 3-4-1 inequality and compact zero-free region.
It is complementary to \texttt{PrimeNumberTheoremAnd}: where the latter uses
modern Fourier analysis, we develop the traditional complex-analytic
machinery that yields explicit error terms.

%=====================================================================
\section{Conclusion and Future Work}
%=====================================================================

We have organized a Lean~4 project that builds against Mathlib~4.29.1 and
collects infrastructure toward de~la~Vall\'ee~Poussin's analytic proof of the
Prime Number Theorem.  The current development is not yet a completed formal
proof of the PNT: it contains no explicit \texttt{sorry} occurrences, but it
still has unproved target statements for quantitative zero-free regions,
Hardy-theorem moment estimates, and explicit-formula machinery.

The immediate path to a publishable artifact is therefore twofold.  First,
keep the repository reproducible and its claims aligned with the checked Lean
source.  Second, prove the project-local target statements before advertising any
zero-free-region or PNT result as \texttt{sorry}-free.  The larger remaining
goals point to useful Mathlib extensions, particularly contour integration,
residue calculus, Borel--Carath\'eodory, and growth estimates for $\zeta$.

Future directions include:
\begin{itemize}[noitemsep]
  \item Completing the Borel--Carath\'eodory route to the classical zero-free
    region $\sigma \geq 1-c/\log|t|$.
  \item Formalizing Perron's formula for the explicit formula.
  \item Extending the framework to Dirichlet $L$-functions for the Prime
    Number Theorem in arithmetic progressions.
\end{itemize}

The full Lean source code is available at
\url{https://github.com/cc-chen-tech/riemann-pnt-lean4}.

%=====================================================================
\section*{Acknowledgments}
This work builds on Mathlib, the Lean~4 community library. We thank the
Mathlib maintainers for the foundational zeta function theory.

%=====================================================================
\bibliographystyle{alpha}
\begin{thebibliography}{99}

\bibitem{avigad2007pnt}
J.~Avigad, K.~Donnelly, D.~Gray, and P.~Raff.
\newblock A formally verified proof of the prime number theorem.
\newblock {\em ACM Trans.\ Comput.\ Logic}, 9(1):2, 2007.

\bibitem{harrison2009pnt}
J.~Harrison.
\newblock Formalizing an analytic proof of the Prime Number Theorem.
\newblock {\em J.\ Autom.\ Reasoning}, 43(3):243--266, 2009.

\bibitem{pntand}
A.~Kontorovich and T.~Tao.
\newblock {PrimeNumberTheoremAnd}.
\newblock \url{https://github.com/AlexKontorovich/PrimeNumberTheoremAnd}, 2024.

\bibitem{mathlib}
The Mathlib Community.
\newblock The Lean mathematical library.
\newblock \url{https://github.com/leanprover-community/mathlib4}, 2024.

\bibitem{davenport}
H.~Davenport.
\newblock {\em Multiplicative Number Theory}, volume~74 of GTM.
\newblock Springer, 3rd edition, 2000.

\bibitem{tenenbaum}
G.~Tenenbaum.
\newblock {\em Introduction to Analytic and Probabilistic Number Theory},
  volume~163 of GSM.
\newblock AMS, 3rd edition, 2015.

\end{thebibliography}

\end{document}
